% Compile with xelatex
\documentclass[11pt,a4paper]{article}

\usepackage{fontspec}
\setmainfont{Malgun Gothic}[Ligatures=TeX]
\setsansfont{Malgun Gothic}[Ligatures=TeX]

\usepackage[top=2.2cm, bottom=2.2cm, left=2.5cm, right=2.5cm]{geometry}
\usepackage{booktabs,longtable,array,multirow}
\usepackage{graphicx}
\usepackage{float}
\usepackage{amsmath}
\usepackage[hidelinks]{hyperref}

\graphicspath{{../data/analysis/figures/}}
\setlength{\parindent}{0pt}
\setlength{\parskip}{6pt}

\newcolumntype{R}{>{\raggedleft\arraybackslash}p{1.4cm}}

% ─────────────────────────────────────────────────────────────
\begin{document}

\begin{center}
{\Large\textbf{Preliminary Analysis Report}}\\[6pt]
{\large Cross-Cultural Perception of Korean Tourist Attractions}\\[4pt]
{\normalsize 8 places (Seoul and Busan) $\cdot$ 472 reviews $\cdot$ October 2026}
\end{center}

\vspace{8pt}
\noindent\rule{\linewidth}{0.4pt}

% ── 1. CORPUS ────────────────────────────────────────────────
\section*{1. Corpus Summary}

A bilingual corpus of 472 reviews was constructed from 8 tourist destinations
in Seoul and Busan. Each place is represented by up to 30 Korean-language
(domestic) and 30 English-language (foreign) reviews, taken in platform order
from culturally segmented online review platforms.

\textbf{Data quality.} The Korean review scraper captured many reviews twice
(once with and once without line breaks, sometimes with a ``video'' badge
prepended). Duplicates are now removed by comparing texts without spaces and
punctuation; in the earlier version of the corpus about a quarter of the
Korean reviews were copies. All results below are computed on the
deduplicated corpus. Only Hongdae Street has fewer than 30 unique Korean
reviews (22).

\begin{table}[H]
\centering
\small
\begin{tabular}{llcc}
\toprule
\textbf{City} & \textbf{Place} & \textbf{KO reviews} & \textbf{EN reviews} \\
\midrule
Seoul  & Gyeongbokgung (경복궁)         & 30 & 30 \\
Seoul  & Changdeokgung (창덕궁)         & 30 & 30 \\
Seoul  & Bukchon Village (북촌한옥마을)   & 30 & 30 \\
Seoul  & Gwangjang Market (광장시장)     & 30 & 30 \\
Seoul  & Hongdae Street (홍대거리)       & 22 & 30 \\
Seoul  & Lotte World (롯데월드)          & 30 & 30 \\
Busan  & Gwangalli Beach (광안리해수욕장) & 30 & 30 \\
Busan  & Haeundae Beach (해운대해수욕장)  & 30 & 30 \\
\midrule
       & \textbf{Total}                 & \textbf{232} & \textbf{240} \\
\bottomrule
\end{tabular}
\caption{Corpus composition by place and language group.}
\end{table}

% ── 2. SENTIMENT ─────────────────────────────────────────────
\section*{2. Sentiment Analysis}

Pre-trained sentiment classifiers were applied to each language group separately.
The table shows the mean sentiment score (range -1 to 1) and the result of a
one-sided Mann-Whitney U test (H$_1$: KO $>$ EN).

\begin{table}[H]
\centering
\small
\begin{tabular}{lRRRRl}
\toprule
\textbf{Place} & \textbf{KO} & \textbf{EN} & \textbf{U} & \textbf{p-value} & \textbf{Sig.} \\
\midrule
Gyeongbokgung    & 0.597 & 0.673 & 397.0 & 0.786 & n.s. \\
Changdeokgung    & 0.589 & 0.668 & 398.0 & 0.781 & n.s. \\
Bukchon Village  & 0.595 & 0.469 & 575.5 & 0.032 & * \\
Gwangjang Market & 0.336 & 0.455 & 460.0 & 0.444 & n.s. \\
Hongdae Street   & 0.666 & 0.419 & 431.0 & 0.031 & * \\
Lotte World      & 0.806 & 0.597 & 569.0 & 0.040 & * \\
Gwangalli Beach  & 0.762 & 0.775 & 364.0 & 0.900 & n.s. \\
Haeundae Beach   & 0.595 & 0.758 & 346.0 & 0.939 & n.s. \\
\midrule
\textbf{Overall} & \textbf{0.617} & \textbf{0.602} & 28796.5 & \textbf{0.259} & n.s. \\
\bottomrule
\end{tabular}
\caption{Mean sentiment and Mann-Whitney U test results. * $p < 0.05$ (uncorrected).}
\end{table}

There is no overall difference between groups ($p = 0.259$). Domestic reviewers
are significantly more positive at three sites: Bukchon Village, Hongdae Street
and Lotte World. Gwangjang Market has the lowest domestic sentiment (0.336).
The three per-site effects do not survive a correction for eight comparisons
and should be read as exploratory.

\begin{figure}[H]
\centering
\includegraphics[width=0.85\linewidth]{01_sentiment_bar.png}
\caption{Mean sentiment by place and language group.}
\end{figure}

% ── 3. CROSS-LINGUAL SIMILARITY ─────────────────────────────
\section*{3. Cross-lingual Semantic Similarity}

Multilingual sentence embeddings (multilingual-e5-base) project Korean and
English reviews into a shared semantic space. Per-site cosine similarity
between group centroids quantifies experiential convergence.

\begin{table}[H]
\centering
\small
\begin{tabular}{lRRR}
\toprule
\textbf{Place} & \textbf{KO-EN sim} & \textbf{KO intra} & \textbf{EN intra} \\
\midrule
Gyeongbokgung    & 0.910 & 0.878 & 0.823 \\
Changdeokgung    & 0.917 & 0.860 & 0.817 \\
Bukchon Village  & 0.898 & 0.878 & 0.813 \\
Gwangjang Market & 0.933 & 0.862 & 0.818 \\
Hongdae Street   & 0.926 & 0.868 & 0.804 \\
Lotte World      & 0.930 & 0.854 & 0.806 \\
Gwangalli Beach  & 0.914 & 0.880 & 0.822 \\
Haeundae Beach   & 0.903 & 0.890 & 0.813 \\
\bottomrule
\end{tabular}
\caption{Cross-lingual cosine similarity and intra-group coherence.}
\end{table}

Cross-group similarity is consistently high (0.90 to 0.93) and exceeds
intra-group similarity at every site, indicating that both groups discuss
similar aspects of each location. Gwangjang Market shows the highest
alignment (0.933); Bukchon Village the lowest (0.898).

% ── 4. MOTIVATION ────────────────────────────────────────────
\section*{4. Travel Motivation (Zero-shot Classification)}

A zero-shot classification using cosine similarity in the embedding space
assigned each review to one of five motivation categories.

\begin{table}[H]
\centering
\small
\begin{tabular}{lRR}
\toprule
\textbf{Motivation} & \textbf{KO (\%)} & \textbf{EN (\%)} \\
\midrule
Photography \& Visual Experience & 36.6 & 37.1 \\
Nature \& Scenery                & 27.2 & 21.7 \\
Food \& Gastronomy               & 15.9 & 15.8 \\
Cultural Heritage \& History     & 14.7 & 15.8 \\
Shopping \& Entertainment        &  5.6 &  9.6 \\
\bottomrule
\end{tabular}
\caption{Overall motivation distribution by language group.}
\end{table}

Both groups are dominated by Photography/Visual ($\sim$37\%). Domestic
reviewers lean more towards Nature \& Scenery (27.2\% vs 21.7\%), while
foreign reviewers show a larger share of Shopping \& Entertainment
(9.6\% vs 5.6\%).

% ── 5. LEXICAL DIVERSITY ────────────────────────────────────
\section*{5. Lexical Diversity (TTR)}

Type-Token Ratio (TTR) measures lexical diversity per review. English reviews
show higher TTR at every site (weighted mean EN = 0.935, KO = 0.871), with the
smallest gap at Lotte World (0.911 vs 0.898). Because TTR depends on how each
language is tokenized, part of this gap may reflect the languages themselves
rather than writing style.

% ── 6. TF-IDF KEYWORDS ──────────────────────────────────────
\section*{6. TF-IDF Keyword Contrast}

\begin{table}[H]
\centering
\small
\begin{tabular}{clcl}
\toprule
\textbf{Rank} & \textbf{Korean} & & \textbf{English} \\
\midrule
1 & 경복궁 (Gyeongbokgung) & & beach \\
2 & 한옥 (hanok)           & & palace \\
3 & 외국인 (foreigners)    & & worth \\
4 & 해운대 (Haeundae)      & & area \\
5 & 시간 (time)            & & experience \\
6 & 북촌 (Bukchon)         & & food \\
7 & 창덕궁 (Changdeokgung) & & people \\
8 & 해수욕장 (beach)       & & street \\
\bottomrule
\end{tabular}
\caption{Top TF-IDF keywords per language group.}
\end{table}

Korean keywords are dominated by place names and by 외국인 (foreigners):
domestic reviewers frequently comment on the presence of foreign visitors.
English keywords are more evaluative (worth, experience).

% ── 7. LDA TOPICS ────────────────────────────────────────────
\section*{7. LDA Topic Modeling}

Five topics were extracted for each language group.

\textbf{English topics:}
\begin{enumerate}\itemsep2pt
  \item beach, clean, enjoy, rides, people, bridge, magic
  \item palace, garden, secret, changdeokgung, traditional, hanbok, history
  \item area, fun, beach, people, food, street, local, shops
  \item beach, busan, walking, haeundae, atmosphere, ocean, sand
  \item beach, food, street, night, area, busan, walk, free
\end{enumerate}

\textbf{Korean topics:}
\begin{enumerate}\itemsep2pt
  \item 한옥, 북촌, 경복궁, 마을, 관광객, 드론, 광안리, 전통, 외국인, 골목
  \item 시간, 경복궁, 풍경, 창덕궁, 경회루, 조선, 사진, 향원정
  \item 해운대, 해수욕장, 창덕궁, 부산, 궁궐, 분위기, 해변, 바다, 축제
  \item 경복궁, 시간, 한복, 사진, 후원, 외국인, 건물, 공연
  \item 홍대, 다양, 풍경, 거리, 한옥, 바다, 외국인, 마을, 부산
\end{enumerate}

\textbf{Choice of K.} Following Chen et al.\ (2024), LDA was fitted for
K = 2 to 15 with three random seeds per K, recording c\_v coherence and
perplexity. Coherence is flat across K with large seed-to-seed variation
(English 0.38 to 0.48; Korean 0.36 to 0.45). The highest mean coherence
is at K = 3 for English (0.480) and K = 8 for Korean (0.447); K = 5 scores
0.425 and 0.400 respectively, within one standard deviation of the best
values. No K is clearly optimal, so K = 5 is retained as a compromise
between interpretability and granularity, and the topics are treated as
descriptive.

\begin{figure}[H]
\centering
\includegraphics[width=\linewidth]{14_lda_k_selection.png}
\caption{Coherence (c\_v) and perplexity by number of topics (mean $\pm$ SD over 3 seeds).}
\end{figure}

% ── 8. HYPOTHESES ────────────────────────────────────────────
\section*{8. Hofstede-based Hypothesis Tests}

Following the design of Chen et al.\ (2024), five hypotheses were derived from
Hofstede's cultural dimensions (Korea: individualism 18, uncertainty avoidance
85, power distance 60, indulgence 29). Each topical hypothesis uses a bilingual
target dictionary. Two tests are run on all reviews of each group:
\textbf{presence} (does a review mention the topic; 2$\times$2 chi-square) and
\textbf{intensity} (share of target tokens per review; Mann-Whitney U). Both are
one-sided and Holm-corrected for five hypotheses. H5 uses the sentiment score
(intensity) and the share of positive reviews, sentiment $> 0.05$ (presence).

\begin{table}[H]
\centering
\small
\begin{tabular}{llrrrrl}
\toprule
\textbf{H} & \textbf{Dimension: topic (expected)} & \textbf{KO \%} & \textbf{EN \%} &
\textbf{OR} & \textbf{p (Holm)} & \textbf{Result} \\
\midrule
H1 & Uncertainty avoidance: planning (KO$>$EN) & 48.3 & 35.0 & 1.73 & 0.009 & Supported \\
H2 & Collectivism: family (KO$>$EN)            & 16.8 &  8.8 & 2.08 & 0.017 & Supported \\
H3 & Power distance: service (KO$>$EN)         &  7.3 &  7.9 & 0.92 & 1.000 & n.s. \\
H4 & Indulgence: enjoyment (EN$>$KO)           & 25.4 & 21.2 & 1.26 & 1.000 & n.s. \\
H5 & Collectivism: positive text (KO$>$EN)     & 87.9 & 85.0 & 1.28 & 0.529 & n.s. \\
\bottomrule
\end{tabular}
\caption{Presence test (chi-square) with full dictionaries. OR = odds ratio KO vs EN.}
\end{table}

In the intensity test only H2 is supported ($p_{Holm} = 0.025$).

\textbf{Dictionary validation.} All 602 matched contexts were coded manually
(keyword in context). Overall precision ranges from 0.73 to 0.97. Six terms
fell below 70\% precision and were removed in a sensitivity run: 시간 (mostly
``a good time''), 버스 (the tokenizer splits 버스킹, busking), 주말, time, times
and friendly.

\begin{table}[H]
\centering
\small
\begin{tabular}{lrrrrl}
\toprule
\textbf{H} & \multicolumn{2}{c}{\textbf{Full KO / EN \%, p}} &
\multicolumn{2}{c}{\textbf{Validated KO / EN \%, p}} & \textbf{Robust} \\
\midrule
H1 Planning & 48.3 / 35.0 & 0.009 & 32.3 / 27.1 & 0.425 & no \\
H2 Family   & 16.8 / 8.8  & 0.017 & 16.8 / 8.8  & 0.022 & yes \\
\bottomrule
\end{tabular}
\caption{Sensitivity of the presence test to dictionary validation (H3 to H5 unchanged, n.s.).}
\end{table}

H1 depended on ambiguous terms and is not robust. \textbf{H2 is the only
robust result}: domestic reviewers are about twice as likely to mention
family or companions (OR = 2.08), consistent with collectivism. The effect
size is small ($\phi = 0.12$).

\textbf{H5 with star ratings.} Text sentiment does not support H5, but star
ratings do in part. Star ratings are available for all English reviews and for
142 of 232 Korean reviews (61\%; coverage varies by place). Among rated reviews,
Korean reviewers rate higher (4.78 vs 4.55; one-sided Mann-Whitney $p = 0.023$),
also after adjusting for place (+0.15 stars, $p = 0.031$). The gap comes from
low ratings: 2.1\% of Korean versus 10.0\% of English ratings are 3 stars or
fewer, while the share of 5-star ratings does not differ after adjustment
(OR = 1.09, n.s.). This matches the wording of Chen et al.\ (2024), that
Chinese-speaking tourists are less inclined to assign low ratings, but rests on a
non-random subset of Korean reviews. As a validation of the text-based measure,
sentiment correlates only weakly with star ratings (Spearman $\rho$ = 0.27 EN,
0.21 KO), partly because almost all ratings are 4 or 5 stars.

\begin{figure}[H]
\centering
\includegraphics[width=\linewidth]{13_hypothesis_tests.png}
\caption{Intensity (top) and presence (bottom) tests per hypothesis, full dictionaries.}
\end{figure}

% ── 9. OBSERVATIONS ──────────────────────────────────────────
\section*{9. Preliminary Observations}

\begin{itemize}\itemsep3pt
  \item No overall sentiment difference between domestic and foreign reviewers
        ($p = 0.259$); differences are site-specific (Bukchon Village,
        Hongdae Street and Lotte World, KO $>$ EN).
  \item High cross-lingual similarity (0.90 to 0.93) shows that both groups
        discuss similar aspects, with different emphasis and vocabulary.
  \item Domestic reviewers mention family and companions about twice as often
        as foreign reviewers; this is the only Hofstede-based hypothesis that
        survives dictionary validation.
  \item Hypotheses on uncertainty avoidance, power distance and indulgence are
        not supported. Rating leniency (H5) is not visible in text sentiment, but
        rated Korean reviews contain fewer low star ratings than English ones.
  \item Limitations: small samples per site, platform and language are
        confounded, and the English-language group is culturally heterogeneous.
  \item Data collection for 2 additional places (인사동길, N서울타워) is pending.
\end{itemize}

\vspace{8pt}
\noindent\rule{\linewidth}{0.4pt}
{\small Hoseo University \hfill October 2026}

\end{document}
